% !TEX program = xelatex
\documentclass{lernblatt}
\thema{Naturschule Chat · Eltern-Anleitung}
\logodatei{logo.png}

\begin{document}
\fancyhead[R]{\scriptsize\textcolor{slate}{Freie Naturschule MSE}}

\station{01}{Windows}{Installation der App „Naturschule Chat“ (Windows 10 \& 11).}

\begin{tipp}{Warum erscheint eine Sicherheitswarnung?}
Als kleine Schule besitzen wir kein kostenpflichtiges Entwickler-Zertifikat von Microsoft oder Apple. Deshalb zeigt Windows beim ersten Start eine Standard-Warnung an. Die App wurde direkt aus unserem Schul-Repository erstellt und ist sicher.
\end{tipp}

\begin{aufg}{Schritt-für-Schritt: Installation unter Windows}
\begin{enumerate}
  \item Lade die Datei \textbf{Naturschule.Chat.msi} herunter und starte sie per Doppelklick.
  \item \textbf{Falls der Browser warnt} (\emph{„Wird nicht häufig heruntergeladen“}): Klicke bei der Meldung auf die drei Punkte (\textbf{\dots}) $\rightarrow$ \textbf{Beibehalten} $\rightarrow$ \textbf{Mehr anzeigen} $\rightarrow$ \textbf{Trotzdem beibehalten}.
  \item \textbf{Windows-SmartScreen-Meldung} (\emph{„Der Computer wurde durch Windows geschützt“}):
    \begin{itemize}
      \item Klicke auf den kleinen Link \textbf{„Weitere Informationen“}.
      \item Danach erscheint unten der Knopf \textbf{„Trotzdem ausführen“} -- klicke darauf.
    \end{itemize}
  \item Folge dem Installationsassistenten bis zum Ende. Danach liegt \textbf{Naturschule Chat} im Startmenü und auf dem Desktop.
\end{enumerate}
\end{aufg}

\begin{falle}{Wichtiger Hinweis für ältere Windows-10-PCs}
Auf älteren Windows-10-Geräten lädt das Installationsprogramm bei Bedarf automatisch die kostenlose Komponente \emph{Microsoft Edge WebView2} nach. Dafür muss während der Installation eine aktive Internetverbindung bestehen.
\end{falle}

\clearpage
\station{02}{macOS (Mac)}{Installation der App „Naturschule Chat“ auf dem Mac.}

\begin{aufg}{Schritt-für-Schritt: Installation auf dem Mac}
\begin{enumerate}
  \item Lade die Datei \textbf{Naturschule.Chat.dmg} herunter und öffne sie per Doppelklick.
  \item Ziehe das Symbol \textbf{Naturschule Chat} in den Ordner \textbf{Programme} (\emph{Applications}).
  \item \textbf{Apple-Gatekeeper-Warnung beim ersten Start:} Da die App nicht von Apple notarisiert ist, blockiert macOS den ersten normalen Doppelklick.
\end{enumerate}
\end{aufg}

\begin{merke}{So öffnest du die App trotzdem (nur beim ersten Mal nötig)}
\textbf{Weg A (Klassisch):}
\begin{itemize}
  \item Öffne im Finder den Ordner \textbf{Programme}.
  \item Mache einen \textbf{Rechtsklick} (oder \textbf{Ctrl + Klick}) auf \textbf{Naturschule Chat} und wähle \textbf{„Öffnen“}.
  \item Bestätige im neuen Fenster erneut mit \textbf{„Öffnen“}.
\end{itemize}

\vspace{2mm}
\textbf{Weg B (Neuere macOS-Versionen wie Sequoia):}
\begin{itemize}
  \item Versuche einmal, die App normal zu starten, und schließe die Warnmeldung mit \textbf{„Fertig“}.
  \item Öffne \textbf{Systemeinstellungen} $\rightarrow$ \textbf{Datenschutz \& Sicherheit} und scrolle ganz nach unten.
  \item Klicke neben \emph{„Naturschule Chat wurde blockiert“} auf \textbf{„Dennoch öffnen“} und bestätige mit deinem Mac-Passwort oder Touch ID.
\end{itemize}
\end{merke}

\station{03}{Erste Schritte}{Login, Benachrichtigungen und Hintergrundbetrieb.}

\begin{tipp}{Damit du keine Chat-Nachrichten verpasst}
\begin{itemize}
  \item \textbf{Einloggen:} Melde dich beim ersten Start einfach mit deinem gewohnten Konto an.
  \item \textbf{Benachrichtigungen erlauben:} Beim ersten Klick in der App erscheint eine Nachfrage für Desktop-Benachrichtigungen -- bitte auf \textbf{Erlauben} klicken.
  \item \textbf{Läuft im Hintergrund weiter:} Wenn du das Fenster schließt (\textbf{X}), bleibt Naturschule Chat in der Symbolleiste (Windows: unten rechts; Mac: oben rechts) aktiv, damit weiterhin Nachrichten ankommen. Zum kompletten Beenden Rechtsklick auf das Symbol $\rightarrow$ \textbf{Quit / Beenden}.
\end{itemize}
\end{tipp}

\end{document}
